The coframe below is that of the derived metric of the Schwarzschild interior gravitational field, the components of $\hb^{-1}$ in an orthonormal frame (Fact \claimref{synthesis.imported-gtg-dictionary}(i)). In units $2M = 1$ the weighted coframe of the Schwarzschild interior carries the weights $r^{1/2}$ on $\sqrt{A}\,dt$ and $r^{-1}$ on $r\,d\theta$, $r\sin\theta\,d\varphi$: these are the Kasner exponents $(-\tfrac13, \tfrac23, \tfrac23)$ of the proper-time variable $\tau \propto r^{3/2}$ rewritten in the variable $r$, where they read $(-\tfrac12, 1, 1)$, normalised to one at $r = r_h$; the horizon-normalised coframe on the exceptional set is therefore $(dt,\ r_h\,d\theta,\ r_h\sin\theta\,d\varphi)$ exactly, $(r^{1/2}\sqrt{A})^2 = 1 - r \to 1$ and $r^{-1}\cdot r = 1$. In the Hawking paper's lift, each channel $(\omega, \ell, m, s)$ of the near-horizon reduction is one two-dimensional mode, and the two helicities of a massless species are two channels at distinct homogeneities $-2s - 2$ ($-4$ and $0$ for the photon, $-6$ and $+2$ for the graviton); the channel identification is that paper's joint J3, its one modelling joint, which it assigns to this paper as an open successor problem, so the lift enters here at modelling tier and not at theorem tier. \emph{Provenance:} \cite[Lemmas 5.1 to 5.3, Def.~5.5]{ButlerT6} (\claimref{interior.interior-connection-exactly}, \claimref{interior.kasner-form-interior}) and \cite[Thm.~8.16]{ButlerT4} (\claimref{sheaf.weights-are-forced-are-kasner-exponents}) for the weights; \cite[Fact~5.1, Prop.~5.2, Assembly~5.3]{ButlerT7} (\claimref{hawking.homogeneity-ordering-shift}, \claimref{hawking.imported-penrose-transform}, \claimref{hawking.helicity-lift-joints-itemised}) for the lift, Assembly 5.3 there itemising J3 as modelling.
